% ===================================================================
%  ТАБЛИЦА № 16 — Большой магазин. — Базар. — Игрушки.
%  Traduction 2026 d'apres texto/io, controlee sur texto/fr, texto/en
%  et texto/es. Consigne : voir texto/ru/10-tablica-01.tex.
%
%  Le tableau d'astronomie (bloc 15-1) porte des renvois a LETTRES,
%  (a) a (f), graves sur le tableau lui-meme : ils se rangent sous le
%  numero 85, comme le dit gravuri/literi.json.
% ===================================================================

%%K t16-apar-1 apar
\VUsaut{4.0mm}
\VUtitre{15.0pt}{60}{ТАБЛИЦА № 16}
\VUsaut{2.0mm}
\VUpk{13.2pt}{40}{Большой магазин. --- Базар.}
\VUpk{13.2pt}{40}{Игрушки.}
\VUsaut{3.4mm}

%%K t16-01-1 p
1. --- Приближается новый год. Большие магазины приготовили свои выставки
\VUgras{игрушек} \textsuperscript{(1)} и новогодних подарков.

%%K t16-01-2 p
Я попросил дедушку сводить меня на базар посмотреть на эту прекрасную
выставку, и он согласился.

%%K t16-02-1 p
2. --- Сперва мы остановились перед красивыми оружейными панно, на которых
были \VUgras{ружьё}, \VUgras{фуражка}, \VUgras{сабля}, \VUgras{портупея} и
пара \VUgras{эполет}, или же \VUgras{каска}, \VUgras{кираса} \textsuperscript{(3)}
и \VUgras{пистолет} \textsuperscript{(4)}. Всё это занимало меня куда больше, чем
\VUgras{волшебные фонари} \textsuperscript{(2)}.

%%K t16-tit-1 sub
\VUsaut{3.0mm}
\VUpk{10.2pt}{40}{Прекрасные зрелища.}
\VUsaut{2.6mm}

%%K t16-03-1 p
3. --- В том же отделе стоял \VUgras{часовой} \textsuperscript{(5)} в своей
\VUgras{будке}; он словно сторожил целый картонный зверинец. Сколько там
было животных! \VUgras{Попугай} \textsuperscript{(6)} на жёрдочке,
\VUgras{носорог} \textsuperscript{(7)} с рогом на носу, \VUgras{северный олень}
\textsuperscript{(8)}, \VUgras{страус} \textsuperscript{(9)}, \VUgras{обезьяна}
\textsuperscript{(10)}, грызущая орех, \VUgras{лисица} \textsuperscript{(11)}, уносящая
курицу, \VUgras{крокодил} \textsuperscript{(12)}, разинувший огромную пасть,
\VUgras{верблюд} \textsuperscript{(13)}, изящная \VUgras{борзая} \textsuperscript{(14)},
\VUgras{пантера} \textsuperscript{(15)}, а потом два зверя, которых я не знал и
которые, как мне сказали, --- \VUgras{ласка} \textsuperscript{(16)} и
\VUgras{гиена} \textsuperscript{(17)}. Был там ещё \VUgras{леопард} \textsuperscript{(18)}
и \VUgras{волк} \textsuperscript{(21)}; совсем рядом --- премилые \VUgras{крысы}
\textsuperscript{(19)} и крошечные резиновые \VUgras{мыши} \textsuperscript{(20)}.
Наконец, \VUgras{бегемот} \textsuperscript{(22)} и горбатый \VUgras{дромадер}
\textsuperscript{(23)} довершали собрание. Я простоял бы там целые часы, если бы
рядом не оказался великолепный лагерь, полный \VUgras{оловянных солдатиков}
\textsuperscript{(29)}, который привлёк моё внимание.

%%K t16-tit-2 sub
\VUsaut{4.0mm}
\VUpk{10.2pt}{40}{Солдаты и война.}
\VUsaut{2.8mm}

%%K t16-04-1 p
4. --- Дедушка, который \VUgras{инвалид} \textsuperscript{(24)} и должен был
оставить службу из-за \VUgras{раны}, принёсшей ему \VUgras{военную медаль},
любит рассказывать мне о \VUgras{походах}, в которых участвовал. Он с
удовольствием объяснил мне разные передвижения этого маленького войска в
миниатюре. Группа настоящих солдат, осматривавших базар --- \VUgras{сапёр}
\textsuperscript{(25)}, \VUgras{альпийский стрелок} \textsuperscript{(26)} в своём
\VUgras{берете}, пехотный \VUgras{фельдфебель} \textsuperscript{(27)} и
артиллерийский \VUgras{сержант} \textsuperscript{(28)} с золотой \VUgras{нашивкой}
на рукаве --- остановилась возле нас послушать объяснения.

%%K t16-05-1 p
5. --- Дедушка, очень гордый таким почтенным слушателем, тут же
сымпровизировал целый план сражения.

%%K t16-05-2 p
Крепость с её \VUgras{валами} и \VUgras{рвами} была осаждена
\VUgras{неприятельским войском}, которое после долгой \VUgras{бомбардировки}
готовилось взять её \VUgras{приступом}. Но \VUgras{осаждённые}, доведённые
голодом, предприняли \VUgras{вылазку} против \VUgras{лагеря}
\VUgras{осаждающих}. Отбив \VUgras{нападение} в жестоком \VUgras{бою} вокруг
\VUgras{палаток}, \VUgras{победившие} осаждающие послали свои
\VUgras{эскадроны} в \VUgras{погоню} за \VUgras{побеждёнными} нападавшими.
Те \VUgras{отступили}, оставив \VUgras{поле битвы} усеянным \VUgras{убитыми}
и \VUgras{ранеными}. \VUgras{Санитары} унесли раненых в
\VUgras{полевой госпиталь}, чтобы их лечил \VUgras{хирург}.

%%K t16-05-3 p
Несмотря на геройское сопротивление нескольких \VUgras{батальонов},
построившихся в \VUgras{каре}, \VUgras{поражение} было полное; остаток
войска \VUgras{бежал}, оставив множество \VUgras{пленных}. Неприятель
готовился увенчать свою \VUgras{победу} взятием крепости. Быть может, то был
бы конец похода, и после \VUgras{перемирия} можно было бы подписать
\VUgras{мирный договор}.

%%K t16-tit-3 sub
\VUsaut{4.4mm}
\VUpk{10.2pt}{40}{Новогодние подарки.}
\VUsaut{2.8mm}

%%K t16-06-1 p
6. --- Молодые солдаты, смеявшиеся, слушая красочный рассказ ветерана,
задали ему ещё несколько вопросов. Что до меня, я обошёл прилавок, чтобы
посмотреть на прекрасный \VUgras{арбалет} \textsuperscript{(32)} и на \VUgras{лук}
\textsuperscript{(33)} с его \VUgras{стрелой} \textsuperscript{(31)}. Там я нашёл ещё
одно собрание животных: двух \VUgras{певчих птиц} \textsuperscript{(34)},
заводных \VUgras{соловья} и \VUgras{славку}, поющих во всё горло, и ряд
странных тварей --- \VUgras{летучую мышь} \textsuperscript{(35)}, \VUgras{удава}
\textsuperscript{(36)}, \VUgras{черепаху} \textsuperscript{(37)}, \VUgras{тюленя}
\textsuperscript{(38)}, \VUgras{кита} \textsuperscript{(39)} и \VUgras{акулу}
\textsuperscript{(40)} из бодрюша.

%%K t16-07-1 p
7. --- Я недолго на них смотрел, потому что заметил то, о чём мечтал:
великолепный \VUgras{кукольный театр} \textsuperscript{(41)} с прекрасными
\VUgras{декорациями} и прелестным красным \VUgras{занавесом}. На
\VUgras{сцене} два актёра словно играли \VUgras{роль} в занимательнейшей
\VUgras{пьесе}, потому что маленькие картонные \VUgras{зрители} в ложах, а
также сидевшие в \VUgras{партере} и на \VUgras{галёрке} позади
\VUgras{дирижёра}, горячо хлопали.

%%K t16-07-2 p
К сожалению, тот театр был очень дорог.

%%K t16-08-1 p
8. --- К тому же выбрать среди игрушек было трудно, потому что витрины были
завалены игрушками всякого рода: \VUgras{Арлекин} \textsuperscript{(42)},
\VUgras{Полишинель} \textsuperscript{(43)}, \VUgras{Пьеро} \textsuperscript{(44)},
\VUgras{совы} \textsuperscript{(45)}, хлопающие крыльями, свистящие
\VUgras{дрозды} \textsuperscript{(46)}, лающие \VUgras{пудели},
\VUgras{граммофон} \textsuperscript{(47)} и \VUgras{головоломки}. Одна из этих
игрушек очень меня позабавила: она изображала \VUgras{морской бой}
\textsuperscript{(48)}, в котором \VUgras{корабль} атакован \VUgras{пиратами} у
берега, засаженного \VUgras{кокосовыми пальмами}.

%%K t16-noto-1 noto
\VUnotes{143.2mm}{%
\textit{(*) См. таблицу № 6.}%
}

%%K t16-09-1 p
9. --- Все эти чудеса я мог разглядывать совсем спокойно, потому что в
магазине было мало народу --- едва горстка зевак, \VUgras{драгун}
\textsuperscript{(49)} и \VUgras{инкассатор} \textsuperscript{(50)}, предъявлявший
\VUgras{вексель} в главной \VUgras{кассе} \textsuperscript{(51)}.

%%K t16-10-1 p
10. --- Я спросил дедушку, для чего книги на полках позади \VUgras{кассира};
он сказал мне, что это \VUgras{конторские книги}.

%%K t16-tit-4 sub
\VUsaut{3.6mm}
\VUpk{10.2pt}{40}{Оглядывая магазин.}
\VUsaut{2.8mm}

%%K t16-11-1 p
11. --- Выйдя из отдела игрушек, мы пошли в шорный, чтобы взглянуть на
\VUgras{сёдла} \textsuperscript{(53)}, \VUgras{стремена} \textsuperscript{(54)},
\VUgras{уздечки} \textsuperscript{(55)}, \VUgras{удила} \textsuperscript{(56)} и
\VUgras{хлысты}. Пока \VUgras{заведующий отделом} \textsuperscript{(57)} давал
дедушке справки, я пошёл смотреть выставку \VUgras{фарфора} и
\VUgras{стекла} \textsuperscript{(58)}, где были прекрасные \VUgras{салатники} и
изящные \VUgras{чайники} \textsuperscript{(59)}.

%%K t16-12-1 p
12. --- Чтобы подняться на первый этаж, мы прошли позади витрины
\VUgras{корсетов} \textsuperscript{(60)}, рядом с чулочным отделом, где были
разложены прекрасные \VUgras{панталоны} \textsuperscript{(63)}. Рядом
\VUgras{приказчик} \textsuperscript{(61)} помогал \VUgras{покупательнице}
\textsuperscript{(62)} выбирать прекрасные шёлковые \VUgras{ткани} \textsuperscript{(64)}.

%%K t16-12-2 p
Поднимаясь по лестнице, я заметил, что магазин украшен японскими
\VUgras{фонариками} \textsuperscript{(65)}, \VUgras{зонтиками} \textsuperscript{(66)} и
огромными \VUgras{пальмами} \textsuperscript{(70)}.

%%K t16-13-1 p
13. --- На первом этаже смотреть было почти нечего: только мебель (*),
\VUgras{несгораемые шкафы} \textsuperscript{(67)}, \VUgras{комоды} \textsuperscript{(68)}
и \VUgras{детские коляски} \textsuperscript{(69)}. Зато общий вид магазина был
очень \VUgras{занятен}.

%%K t16-14-1 p
14. --- Под нами я видел музыкальные инструменты: \VUgras{арфы}
\textsuperscript{(71)}, \VUgras{гитары} \textsuperscript{(72)}, \VUgras{мандолины}
\textsuperscript{(73)}, \VUgras{корнеты} \textsuperscript{(74)}, \VUgras{кларнеты}
\textsuperscript{(75)}, скрипки с их \VUgras{смычками} \textsuperscript{(76)} и даже
\VUgras{большой барабан} \textsuperscript{(77)}. Немного дальше, у писчебумажного
прилавка, \VUgras{студент} \textsuperscript{(78)} торговался с
\VUgras{приказчиком} \textsuperscript{(79)} о папке для бумаг. Возле них я
различал прекрасную \VUgras{азбуку} \textsuperscript{(80)}.

%%K t16-14-2 p
Посреди магазина стояла высокая \VUgras{рождественская ёлка} \textsuperscript{(81)},
вся увешанная свечами, безделушками и игрушками всякого рода:
\VUgras{деревянными лошадками} \textsuperscript{(82)}, \VUgras{куклами}
\textsuperscript{(83)} и прочим.

%%K t16-14-3 p
\VUgras{Парфюмерный отдел} \textsuperscript{(84)} со своими \VUgras{зубными
эликсирами} и разными \VUgras{духами} издавал очень приятный запах, который
поднимался к нам.

%%K t16-tit-5 sub
\VUsaut{3.4mm}
\VUpk{10.2pt}{40}{Мы кое-что покупаем.}
\VUsaut{2.8mm}

%%K t16-15-1 p
15. --- Так как дедушке становилось скучно, мы снова спустились по большой
лестнице. Он на минуту остановился перед \VUgras{астрономической таблицей}
\textsuperscript{(85)}, на которой, как он мне сказал, были изображены
\VUgras{кометы} \textsuperscript{(a)}, \VUgras{затмения луны и солнца}
\textsuperscript{(b)}, роза ветров с \VUgras{четырьмя странами света}
\textsuperscript{(c)} \VUgras{(север, юг, восток и запад)}, карта двух
\VUgras{полушарий} \textsuperscript{(d)}, \VUgras{планеты} \textsuperscript{(e)} и
\VUgras{солнечная система} \textsuperscript{(f)}. Я ничего в этом не понял, и
меня куда больше занимала галунная ливрея проходившего мимо
\VUgras{рассыльного} \textsuperscript{(86)} и хорошенькие \VUgras{веера}
\textsuperscript{(87)} рядом со мной, возле \VUgras{кошельков} \textsuperscript{(88)} и
\VUgras{бумажников} \textsuperscript{(89)}.

%%K t16-16-1 p
16. --- Я любовался также на прилавке \VUgras{перьями} \textsuperscript{(90)} и
\VUgras{пряжками для поясов} \textsuperscript{(91)}, и хорошенькими
\VUgras{искусственными цветами} \textsuperscript{(94)}, изящно уложенными в
корзинах. \VUgras{Фиалки}, \VUgras{левкои}, \VUgras{гиацинты} \textsuperscript{(95)},
\VUgras{герани} \textsuperscript{(96)}, \VUgras{лилии} \textsuperscript{(97)} и
\VUgras{настурции} \textsuperscript{(98)} были очень хорошо подделаны.

%%K t16-tit-6 sub
\VUsaut{4.4mm}
\VUpk{10.2pt}{40}{Дедушкин подарок.}
\VUsaut{2.8mm}

%%K t16-17-1 p
17. --- Я попросил дедушку купить мне одну из \VUgras{зубных щёток}
\textsuperscript{(92)} в коробке возле \VUgras{шпилек} \textsuperscript{(93)}. Он
согласился, и я сам пошёл заплатить за свою покупку в кассу, возле которой
увязывали большой \VUgras{свёрток} \textsuperscript{(100)} для одной
\VUgras{покупательницы} \textsuperscript{(99)}.

%%K t16-18-1 p
18. --- Вдруг среди людей, стоявших возле художественных \VUgras{бронз}
\textsuperscript{(101)}, поднялось лёгкое смятение. \VUgras{Вора} \textsuperscript{(102)}
только что поймал с поличным \VUgras{магазинный сыщик} \textsuperscript{(103)} в
ту минуту, когда тот пытался кого-то обокрасть. Послали за городовым, чтобы
его \VUgras{задержать}, и как раз когда мы выходили, виновного вели в
участок.
\VUsaut{6.0mm}
\VUfilet{22mm}
